%%
%    * ----------------------------------------------------------------
%    * "THE BEER-WARE LICENSE" (Revision 42/023):
%    * Ronny Bergmann <mail@rbergmann.info> wrote this file. As long as
%    * you retain this notice you can do whatever you want with this
%    * stuff. If we meet some day and you think this stuff is worth it,
%    * you can buy me a beer or a coffee in return.
%    * ----------------------------------------------------------------
%
%
% A german example using the Kartei.cls - including print and toc as
% options, hence all pages are Din A4.
%
% Last Change: Kartei 1.9, 2012/01/04
%
\documentclass[a9paper,9pt,grid=front%
,toc
%,print
]{kartei}
%\usepackage[utf8]{inputenc} %UTF8
%UTF8
\usepackage{xltxtra}

\usepackage[xetex]{hyperref}

\begin{document}
\setcardpagelayout
\begin{karte}[Notdünger]{
\includegraphics{../regular/bulbasaur}
\includegraphics{../shiny/bulbasaur} \\
Saurier; 
Samen
}[Chlorophyll]
Bisasam: Pflanze, Gift.

r Saurier; -s, -

r Samen; -s, -

\end{karte}

\begin{karte}[Zahlenkunde]{Was ist der Unterschied in der Verwendung von Drölf 
und $n$ bei Ihnen?}
$n$ wird verwendet für Zahlen bis hin zu verdammt groß, Drölf nur bis hin zu 
verdammt.

\end{karte}
%  \section{Informatik}
%  \subsection{Spaß mit Verweisen}
%  %Den Kommentar im Stil ändern

\renewcommand{\kommentarstil}{\textsc}

\begin{karte}{Was ist verschränkte Rekursion ?}[ein Beispiel für Label]
\label{karte:antwort} Siehe Karte \ref{karte:frage}
\end{karte}
%  \subsection{Spaß mit Verweisen II}
\begin{karte}{Was ist die Antwort auf Karte \ref{karte:antwort} ?}
\label{karte:frage}
  Hier kommt man eigentlich gar nicht hin. Hier gibt es also nichts zu sehen, 
bitte blättern sie unauffällig weiter.
\end{karte}

\end{document}